\documentclass[11pt]{article}
\usepackage[T1]{fontenc}
\usepackage{lmodern}
\usepackage[margin=1in]{geometry}
\usepackage{amsmath,amssymb,amsthm}
\usepackage{graphicx,tikz-cd}
\usepackage[hypertexnames=false,colorlinks=true,linkcolor=blue,urlcolor=blue]{hyperref}
\newcommand{\Z}{\mathbb Z}
\newcommand{\F}{\mathcal F}
\newcommand{\G}{\mathcal G}
\newcommand{\shHom}{\underline{\operatorname{Hom}}}
\newcommand{\shExt}{\underline{\operatorname{Ext}}}
\newcommand{\Ext}{\operatorname{Ext}}
\newcommand{\Hom}{\operatorname{Hom}}
\newcommand{\Sh}{\operatorname{Sh}}
\newcommand{\Pic}{\operatorname{Pic}}
\newcommand{\Div}{\operatorname{Div}}
\newcommand{\Jac}{\operatorname{Jac}}
\newcommand{\et}{\mathrm{et}}
\newcommand{\loc}[1]{\par\smallskip\noindent{\small\textit{Source: #1}}\par\smallskip}
\newcommand{\step}[1]{\par\medskip\noindent\textit{#1}\quad}
\newsavebox{\diagramBox}
\newcommand{\fitdiagram}[1]{\par\medskip\noindent\sbox{\diagramBox}{#1}\ifdim\wd\diagramBox>\linewidth\resizebox{\linewidth}{!}{\usebox{\diagramBox}}\else\makebox[\linewidth]{\usebox{\diagramBox}}\fi\par\medskip}
\setlength{\emergencystretch}{2em}
\begin{document}
\begin{center}
{\Large Poincar\'e duality for a regular curve}\par\medskip
{\large J.S. Milne, \textit{Lectures on Etale Cohomology}}\par
Theorem 14.7 and its printed proof sketch\par\smallskip
Version 2.21, March 22, 2013
\end{center}
\noindent\textbf{Private reading draft.} This is a checked-source transcription draft of selected author material from the fixed 202-page PDF (1,514,442 bytes), not a new proof or a mathematical correctness certification. The primary statement and complete printed sketch are on source pp.\ 93--96. The deferred construction of the specific pairing on pp.\ 102--103 is supplied as uncounted context. Earlier strong inputs are quoted as statements with their original locators; their proofs are not all reprinted here. Source: \url{https://www.jmilne.org/math/CourseNotes/LEC.pdf}. Public redistribution of an edited reader has not been established.

% BEGIN CONTEXT LEC:14.7
\section*{Selected source context (not separate targets)}
\noindent\textit{Editorial scope.} The source's curve, sheaf, and finite-coefficient conventions below are retained. In particular, no properness or characteristic-zero restriction is added. The isolated statements collected here are inputs to the printed proof, rather than additional primary results. Source section references in the transcription refer to Milne's notes.

\subsection*{Sheaves, finite direct image, and extension by zero}
\loc{Proposition 6.16, p.\ 49.}
Let $X$ be connected, and let $\bar x$ be a geometric point of $X$. The functor $\F\mapsto\F_{\bar x}$ defines an equivalence between the category of locally constant sheaves of sets (resp.\ abelian groups) with finite stalks and the category of finite $\pi_1(X,\bar x)$-sets (resp.\ modules).

\loc{Corollary 8.4, p.\ 59; recalled in Step 1, p.\ 94.}
If $\pi\colon Y\to X$ is finite or a closed immersion, then $\pi_*$ is exact. In Step 1 the source also recalls that $\pi_*$ preserves injectives. These are source inputs; no proof is added here. The finite-direct-image stalk statement is Proposition 8.3(c), whose proof the source omits.

\loc{Extension by zero and Propositions 8.13--8.15, p.\ 62.}
Let $j\colon U\hookrightarrow X$ be an open immersion. For a presheaf $P$ on $U$, define $P_!(V)=P(V)$ if the image of $\varphi\colon V\to X$ is contained in $U$, and $P_!(V)=0$ otherwise. For a sheaf $\F$ on $U$, $j_!\F$ is the sheaf associated with $\F_!$. The functor $j_!$ is left adjoint to $j^*$:
\[
\Hom_X(j_!\F,\G)\simeq\Hom_U(\F,\G|U).
\]
Its stalk at $\bar x$ is $\F_{\bar x}$ if $x\in U$, and zero otherwise. It is exact, and $j^*$ preserves injectives. If $i\colon Z\hookrightarrow X$ is the complementary closed immersion, then
\[
0\longrightarrow j_!j^*\F\longrightarrow\F\longrightarrow i_*i^*\F\longrightarrow0
\]
is exact. The printed Step 2 retains its original reference to (8.16).

\loc{Example 11.3, p.\ 78, as used with $G=\Z/n\Z$.}
Let $G$ be a constant finite group sheaf on $X_{\et}$. A Galois covering of $X$ with group $G$ is a principal homogeneous space for $G$, and every principal homogeneous space arises in this way. When $X$ is connected, the Galois coverings of $X$ with group $G$ are classified by the continuous homomorphisms $\pi_1(X,\bar x)\to G$. Thus, in the notation of the source,
\[
H^1(X_{\et},G)\simeq\Hom_{\mathrm{conts}}(\pi_1(X,\bar x),G).
\]

\subsection*{The curve inputs}
\loc{Proposition 14.1 and its preceding definitions, pp.\ 89--90.}
Let $X$ be a complete connected nonsingular curve over an algebraically closed field. The divisor $\operatorname{div}(f)$ of a nonzero rational function $f$ on $X$ has degree zero. Let $\Div^0(X)$ denote the group of divisors of degree zero, and $\Pic^0(X)$ the quotient of $\Div^0(X)$ by the subgroup of principal divisors. The sequence
\[
0\longrightarrow\Pic^0(X)\longrightarrow\Pic(X)\longrightarrow\Z\longrightarrow0
\]
is exact. For any integer $n$ relatively prime to the characteristic of $k$,
\[
z\longmapsto nz\colon\Pic^0(X)\longrightarrow\Pic^0(X)
\]
is surjective with kernel equal to a free $\Z/n\Z$-module of rank $2g$, where $g$ is the genus of $X$.

\loc{Proposition 14.2, p.\ 90.}
Let $X$ be a complete connected nonsingular curve over an algebraically closed field $k$. For any $n$ prime to the characteristic of $k$,
\[
H^0(X,\mu_n)=\mu_n(k),\qquad
H^1(X,\mu_n)\approx(\Z/n\Z)^{2g},\qquad
H^2(X,\mu_n)\simeq\Z/n\Z,
\]
and $H^r(X,\mu_n)=0$, $r>2$.

\loc{Proposition 14.3, p.\ 91.}
Let $U$ be a nonsingular curve over an algebraically closed field $k$. For any $n$ prime to the characteristic of $k$ and closed point $x\in U$,
\[
H_x^2(U,\mu_n)\simeq\Z/n\Z,\qquad H_x^r(U,\mu_n)=0\quad\text{for }r\ne2.
\]

\loc{Remark 14.4, p.\ 91.}
Let $M$ be a free $\Z/n\Z$-module of rank $1$, and let $M$ denote also the constant sheaf on a variety (or scheme) $Y$ defined by $M$. Then
\[
H^r(Y_{\et},M)\simeq H^r(Y_{\et},\Z/n\Z)\otimes M
\approx H^r(Y_{\et},\Z/n\Z).
\]
This remark applies to $M=\mu_n$ when $Y$ is a variety over a field $k$ containing $n$ distinct roots of $1$.

\subsection*{Completion and compact support}
\loc{Cohomology with compact support for curves, pp.\ 91--92.}
A function field in one variable over a field $k$ is a field $K\supset k$ of transcendence degree $1$ over $k$ and finitely generated over $k$, i.e., $K$ is a finite extension of $k(T)$ for some $T$ transcendental over $k$. A curve $U$ over $k$ is regular if $\mathcal O_{U,u}$ is a discrete valuation ring for all $u$.

Every function field $K$ in one variable has a connected complete regular curve $X$ canonically associated with it. In fact, we can define $X$ to be the set of all discrete valuation rings $R$ in $K$ satisfying
\[
(*)\qquad R\text{ contains }k\text{ and has field of fractions }K,
\]
and endow it with the topology for which the proper closed subsets are the finite sets. For each open $U$ in $X$, define
\[
\Gamma(U,\mathcal O_X)=\bigcap_{R\in U}R.
\]
Then $(X,\mathcal O_X)$ is a complete curve with $K$ as its field of rational functions, and the local ring at the point $R$ of $X$ is $R$ itself. (See Hartshorne 1977, I.6.)

Let $U$ be a connected regular curve over $k$, and let $X$ be the connected complete regular curve canonically associated with $k(U)$. Then the map
\[
j\colon U\longrightarrow X,\qquad j(u)=\mathcal O_{U,u},
\]
realizes $U$ as an open subvariety of $X$. For a sheaf $\F$ on $U$, we define
\[
H_c^r(U,\F)=H^r(X,j_!\F)
\]
and refer to the $H_c^r(U,\F)$ as the cohomology groups of $\F$ with compact support. For an explanation of this terminology, see \S18 below.

Because $j_!$ is exact, a short exact sequence of sheaves on $U$ gives a long exact sequence of cohomology groups. However, because $j_!$ doesn't preserve injectives, $H_c^r(U,\F)$ is not the $r$th right derived functor of $H_c^0(U,-)$.

\loc{Proposition 14.5, p.\ 92.}
For any connected regular curve $U$ over an algebraically closed field $k$ and integer $n$ not divisible by the characteristic of $k$, there is a canonical isomorphism
\[
H_c^2(U,\mu_n)\longrightarrow\Z/n\Z.
\]

\loc{Remark 14.6, p.\ 92.}
For any $x\in U$, sheaf $\F$ on $U$, and $r\geq0$, there is a canonical map $H_x^r(U,\F)\to H_c^r(U,\F)$. For $\F=\mu_n$ and $r=2$, the map is compatible with the isomorphisms in (14.3) and (14.5).

\subsection*{The vanishing input used in Step 0}
\loc{Theorem 15.1, first statement only, p.\ 105; its printed proof is on pp.\ 105--106.}
A variety (or scheme) $X$ is said to have cohomological dimension $c$ if $c$ is the least integer (or $\infty$) such that $H^r(X_{\et},\F)=0$ for $r>c$ and all torsion sheaves $\F$ on $X_{\et}$. For a variety $X$ over an algebraically closed field $k$,
\[
\operatorname{cd}(X)\leq2\dim(X).
\]
\noindent\textit{Input boundary.} This is a supplied same-book strong input, not an assumption of duality itself. Its first-statement proof on pp.\ 105--106 is present in the fixed source and material input, but is not reprinted in this reader. It uses Theorem 15.2, for which the author refers to Shatz 1972, p.\ 119, Theorem 28. The affine bound and its separate proof on p.\ 107 are not used here. Similarly, the earlier proof of Proposition 14.3 has its original dependencies, including the source's prior omission in (13.6)(c); this reader does not supply a new proof of that input.
% END CONTEXT LEC:14.7

\clearpage
% BEGIN PRIMARY STATEMENT LEC:14.7
\section*{The Poincar\'e duality theorem}
\loc{Theorem 14.7, definitions, and comments, p.\ 93.}
Throughout this subsection, $U$ is a connected regular curve over an algebraically closed field $k$, and $n$ is an integer not divisible by the characteristic of $k$.

By a finite locally constant sheaf $\F$, I shall mean a locally constant sheaf, killed by $n$, that has finite stalks. Thus, for some finite \'etale covering $U'\to U$, $\F|U'$ is the constant sheaf defined by a finite $\Z/n\Z$-module $M$, and to give a finite locally constant sheaf $\F$ on $U$ is to give a finite $\Z/n\Z$-module endowed with a continuous action of $\pi_1(U,\bar u)$. (See \S6.)

\medskip\noindent\textbf{Theorem 14.7 (Poincar\'e Duality).}
\textit{For any finite locally constant sheaf $\F$ on $U$ and integer $r\geq0$, there is a canonical perfect pairing of finite groups}
\[
H_c^r(U,\F)\times H^{2-r}(U,\check{\F}(1))
\longrightarrow H_c^2(U,\mu_n)\simeq\Z/n\Z.
\]
A pairing $M\times N\to C$ is said to be perfect if the induced maps
\[
M\longrightarrow\Hom(N,C),\qquad N\longrightarrow\Hom(M,C)
\]
are isomorphisms. The sheaf $\check{\F}(1)$ is
\[
V\longmapsto\Hom_V(\F|V,\mu_n|V).
\]
Note that, if $\G=\check{\F}(1)$, then $\F=\check{\G}(1)$.

For a discussion of the relation of the above theorem to the usual Poincar\'e duality theorem for topological spaces, see \S24.

I defer a description of the pairing until later. In the sketch of the proof of the theorem, I omit proofs that the diagrams commute. It is important to understand though, that the theorem is about a specific pairing, i.e., it doesn't just say that the groups are dual, but that they are dual with a specific pairing.
% END PRIMARY STATEMENT LEC:14.7

\medskip\noindent\textbf{Example 14.8 (uncounted adjacent comment, p.\ 93).}
Let $X$ be a complete connected nonsingular curve over an algebraically closed field $k$. In this case, the pairing is
\[
H^r(X,\F)\times H^{2-r}(X,\check{\F}(1))
\longrightarrow H^2(X,\mu_n)\simeq\Z/n\Z.
\]
On taking $\F=\Z/n\Z$ and tensoring two of the groups with $\mu_n$ (cf.\ 14.4), we obtain a pairing
\[
H^r(X,\mu_n)\times H^{2-r}(X,\mu_n)
\longrightarrow H^2(X,\mu_n\otimes\mu_n)\simeq\mu_n.
\]
For $r=1$, this can be identified with the Weil pairing
\[
\Jac(X)_n\times\Jac(X)_n\longrightarrow\mu_n
\]
on the points of order $n$ on the Jacobian of $X$ (AV I \S13).

\clearpage
% BEGIN PRIMARY PROOF LEC:14.7
\section*{Sketch of the proof of Poincar\'e duality}
\loc{Complete printed sketch, pp.\ 93--96.}
In the proof, we make frequent use of the five-lemma in the following form. Consider a commutative diagram with exact rows,
\fitdiagram{$\begin{tikzcd}[ampersand replacement=\&,column sep=3em,row sep=2em]
\bullet\arrow[r]\arrow[d,"a"']\&\bullet\arrow[r]\arrow[d,"b"']\&\bullet\arrow[r]\arrow[d,"c"']\&\bullet\arrow[d,"d"]\\
\bullet\arrow[r]\&\bullet\arrow[r]\&\bullet\arrow[r]\&\bullet
\end{tikzcd}$}
then
\[
\begin{aligned}
b,d\text{ injective},\ a\text{ surjective}&\ \Longrightarrow\ c\text{ injective};\\
a,c\text{ surjective},\ d\text{ injective}&\ \Longrightarrow\ b\text{ surjective}.
\end{aligned}
\]

\step{Step 0.} For any torsion sheaf $\F$, the groups $H^r(U,\F)$ and $H_c^r(U,\F)$ are zero for $r>2$.
\begin{proof}
For the vanishing of $H^r(U,\F)=0$, see Theorem 15.1 below. It implies that $H_c^r(U,\F)\overset{\mathrm{def}}=H^r(X,j_!\F)=0$ for $r>2$.
\end{proof}
Thus the theorem is true for $r\ne0,1,2$.

For convenience, write $T^r(U,\F)=H_c^{2-r}(U,\check{\F}(1))^\vee$ (the last $\vee$ means dual in the sense $\Hom(-,\Z/n\Z)$). We have to show that the map
\[
\phi^r(U,\F)\colon H^r(U,\F)\longrightarrow T^r(U,\F)
\]
defined by the pairing in the theorem (with $\F$ replaced by $\check{\F}(1)$) is an isomorphism of finite groups for all finite locally constant sheaves $\F$ on $U$.

Because $\Hom_{\Z/n\Z}(-,\Z/n\Z)$ preserves the exactness of sequences of $\Z/n\Z$-modules, so also does $\F\mapsto\check{\F}(1)$ (look on stalks). It follows that an exact sequence of finite locally constant sheaves
\[
0\longrightarrow\F'\longrightarrow\F\longrightarrow\F''\longrightarrow0
\]
gives rise to an exact sequence
\[
\cdots\longrightarrow T^s(U,\F')\longrightarrow T^s(U,\F)
\longrightarrow T^s(U,\F'')\longrightarrow T^{s+1}(U,\F')\longrightarrow\cdots.
\]

\step{Step 1:} Let $\pi\colon U'\to U$ be a finite map. The theorem is true for $\F$ on $U'$ if and only if it is true for $\pi_*\F$ on $U$.
\begin{proof}
Recall that $\pi_*$ is exact and preserves injectives, and so $H^r(U,\pi_*\F)\simeq H^r(U',\F)$ for all $r$. Since a similar statement is true for $T^s$, we see that $\phi^r(U,\pi_*\F)$ can be identified with $\phi^r(U',\F)$.
\end{proof}

\step{Step 2:} Let $V=U\smallsetminus x$ for some point $x\in U$. Then there is an exact commutative diagram with isomorphisms where indicated:
\fitdiagram{$\begin{tikzcd}[ampersand replacement=\&,column sep=1.05em,row sep=3.2em]
H_x^r(U,\mu_n)\arrow[r]\arrow[d,"\approx"']\&H^r(U,\mu_n)\arrow[r]\arrow[d,"{\phi^r(U,\mu_n)}"']\&H^r(V,\mu_n)\arrow[r]\arrow[d,"{\phi^r(V,\mu_n)}"']\&H_x^{r+1}(U,\mu_n)\arrow[d,"\approx"]\\
H^{2-r}(x,\Z/n\Z)^\vee\arrow[r]\&H_c^{2-r}(U,\Z/n\Z)^\vee\arrow[r]\&H_c^{2-r}(V,\Z/n\Z)^\vee\arrow[r]\&H^{3-r}(x,\Z/n\Z)^\vee
\end{tikzcd}$}
\begin{proof}
The upper sequence is the exact sequence of the pair $(U,V)$. According to (14.3), $H_x^r(V,\mu_n)=H^0(x,\Z/n\Z)$ when $r=2$ and is zero otherwise.

The lower sequence is the compact cohomology sequence of
\[
0\longrightarrow j_!\Z/n\Z\longrightarrow\Z/n\Z
\longrightarrow i_*i^*\Z/n\Z\longrightarrow0,
\]
where $j$ and $i$ are the inclusions of $V$ and $x$ into $U$ respectively (see 8.16).

For $r=2$, the unlabelled vertical map is the isomorphism given by the pairing
\[
H^0(x,\Z/n\Z)\times H^0(x,\Z/n\Z)
\longrightarrow H^0(x,\Z/n\Z)\simeq\Z/n\Z
\]
(Poincar\'e duality for the point $x$), and for $r\ne2$, it is a map between zero groups.
\end{proof}

\step{Step 3:} The map $\phi^0(U,\Z/n\Z)$ is an isomorphism of finite groups.
\begin{proof}
In this case, the pairing is
\[
H^0(U,\Z/n\Z)\times H_c^2(U,\mu_n)
\longrightarrow H_c^2(U,\mu_n)\simeq\Z/n\Z.
\]
Here $H^0(U,\Z/n\Z)=\Z/n\Z$, and its action on $H_c^2(U,\mu_n)$ is defined by the natural $\Z/n\Z$-module structure on $H_c^2(U,\mu_n)$.
\end{proof}

\step{Step 4:} The theorem is true for $r=0$ and $\F$ locally constant.
\begin{proof}
Let $U'\to U$ be a finite \'etale covering such that $\F|U'$ is constant. We can embed $\F|U'$ into a sheaf $\F'=(\Z/n\Z)^s$ on $U'$ for some $s$. On applying $\pi_*$ to $\F|U'\hookrightarrow\F'$ and composing the result with the natural inclusion $\F\hookrightarrow\pi_*\pi^*\F$, we obtain the first map in the sequence
\[
0\longrightarrow\F\longrightarrow\pi_*\F'\longrightarrow\F''\longrightarrow0.
\]
The cokernel $\F''$ is again locally constant (when $U'\to U$ is chosen to be Galois, $\F''$ becomes constant on $U'$). Consider the diagram
\fitdiagram{$\begin{tikzcd}[ampersand replacement=\&,column sep=1.7em,row sep=2.3em]
\bullet\arrow[r]\arrow[d,"\approx"']\&\bullet\arrow[r]\arrow[d,"\approx"']\&H^0(U,\F)\arrow[r]\arrow[d]\&H^0(U,\pi_*\F')\arrow[r]\arrow[d,"\approx"']\&H^0(U,\F'')\arrow[d]\\
\bullet\arrow[r]\&\bullet\arrow[r]\&T^0(U,\F)\arrow[r]\&T^0(U,\pi_*\F')\arrow[r]\&T^0(U,\F'')
\end{tikzcd}$}
The groups ``$\bullet$'' are in fact zero, but all we need is that the arrows are isomorphisms. The map $\phi^0(U,\pi_*\F')$ is an isomorphism by Steps 1 and 3. The five-lemma shows that $\phi^0(U,\F)$ is a injective. Since this is true for all locally constant sheaves $\F$, $\phi^0(U,\F'')$ is injective, and the five-lemma now shows that $\phi^0(U,\F)$ is surjective. Finally, $H^0(U,\F)$ is obviously finite.
\end{proof}

\step{Step 5:} The map $\phi^1(U,\mu_n)$ is injective.
\begin{proof}
Recall (11.3) that $H^1(U,\Z/n\Z)=\Hom_{\mathrm{conts}}(\pi_1(U,\bar u),\Z/n\Z)$. Let $s\in H^1(U,\Z/n\Z)$, and let $\pi\colon U'\to U$ be the Galois covering corresponding to the kernel of $s$. Then $s$ maps to zero in $H^1(U',\Z/n\Z)$, which is isomorphic to $H^1(U,\pi_*\Z/n\Z)$. Let $\F''$ be the cokernel of $\Z/n\Z\to\pi_*(\Z/n\Z)$, and consider
\fitdiagram{$\begin{tikzcd}[ampersand replacement=\&,column sep=1.4em,row sep=2.5em]
H^0(U,\pi_*\Z/n\Z)\arrow[r]\arrow[d,"\approx"']\&H^0(U,\F'')\arrow[r]\arrow[d,"\approx"']\&H^1(U,\Z/n\Z)\arrow[r]\arrow[d]\&H^1(U,\pi_*\Z/n\Z)\arrow[d]\\
T^0(U,\pi_*\Z/n\Z)\arrow[r]\&T^0(U,\F'')\arrow[r]\&T^1(U,\Z/n\Z)\arrow[r]\&T^1(U,\pi_*\Z/n\Z)
\end{tikzcd}$}
From our choice of $\pi$, $s$ maps to zero in $H^1(U,\pi_*\Z/n\Z)$, and a diagram chase shows that, if $s$ also maps to zero in $T^1(U,\Z/n\Z)$, then it is zero.
\end{proof}

\step{Step 6.} The maps $\phi^r(U,\Z/n\Z)$ are isomorphisms of finite groups for $r=1,2$.
\begin{proof}
First take $U=X$. For $r=1$, $\phi$ is a map $H^1(X,\Z/n\Z)\to H^1(X,\mu_n)^\vee$. Because $\phi^1(X,\Z/n\Z)$ is injective (Step 5), and $H^1(X,\Z/n\Z)$ is finite of the same order as $H^1(X,\mu_n)$ (14.2, 14.4), it is an isomorphism. For $r=2$, we have to show that the pairing
\[
H^2(X,\Z/n\Z)\times\mu_n(k)\longrightarrow H^2(X,\mu_n)
\]
is perfect, but this follows from (14.4).

We have shown that $\phi^r(X,\Z/n\Z)$ is an isomorphism of finite groups for $r=1,2$. To deduce the same statement for $\phi^r(U,\Z/n\Z)$, remove the points of $X\smallsetminus U$ one at a time, and apply the five-lemma to the diagram in Step 2.
\end{proof}

\step{Step 7.} The maps $\phi^r(U,\F)$ are isomorphisms of finite groups for $r=1,2$ and $\F$ locally constant.
\begin{proof}
Apply the argument in Step 4, twice.
\end{proof}
This completes the sketch of the proof of Theorem 14.7.
% END PRIMARY PROOF LEC:14.7

\clearpage
% BEGIN PAIRING CONTEXT LEC:14.7
\section*{The deferred specific pairing (uncounted context)}
\loc{The Ext group; Lemma 14.21 and its complete printed sketch, p.\ 102.}
For any variety (or scheme) $X$, let $\Sh(X_{\et},n)$ be the category of sheaves of $\Z/n\Z$-modules on $X_{\et}$. Just as for the full category, $\Sh(X_{\et},n)$ has enough injectives. For a sheaf $\F$ on $\Sh(X_{\et},n)$, $H^r(X_{\et},\F)$ will be the same whether computed using injective resolutions in $\Sh(X_{\et})$ or in $\Sh(X_{\et},n)$ (see EC III 2.25). However, if $\F$ and $\G$ are two sheaves killed by $n$, then $\Ext^r(\F,\G)$ will depend on which category we compute it in. This can be seen already when $X$ is the spectrum of an algebraically closed field. Here $\Sh(X_{\et})$ is the category of $\Z$-modules, and $\Sh(X_{\et},n)$ is the category of $\Z/n\Z$-modules. In the first category, $\Ext^1(\Z/n\Z,\Z/n\Z)\ne0$, because
\[
0\longrightarrow\Z/n\Z\longrightarrow\Z/n^2\Z
\xrightarrow{n}\Z/n\Z\longrightarrow0
\]
doesn't split, whereas in the second category $\Ext^1(\Z/n\Z,\Z/n\Z)=0$ because $\Z/n\Z$ is injective (as a $\Z/n\Z$-module).

By $\Ext^r_{U,n}(\F,\G)$ we mean the Ext group computed in the category $\Sh(U_{\et},n)$.

\medskip\noindent\textbf{Lemma 14.21.}
\textit{If $\F$ is locally constant, then}
\[
\Ext^r_{U,n}(\F,\mu_n)\simeq H^r(U_{\et},\check{\F}(1))
\]
\textit{where $\check{\F}(1)$ is the sheaf}
\[
V\longmapsto\Hom_V(\F|V,\mu_n).
\]
\begin{proof}[Proof (Sketch of Proof)]
For sheaves $\F$ and $\G$ on an arbitrary variety (or scheme) $Y$, we define $\shHom(\F,\G)$ to be the sheaf
\[
V\longmapsto\Hom_{V_{\et}}(\F|V,\G|V)
\]
--- it is easy to check that $\shHom(\F,\G)$ is in fact a sheaf. Let $\F_0$ be a locally constant sheaf of $\Z/n\Z$-modules with finite stalks. The functor
\[
\G\longmapsto\shHom(\F_0,\G)\colon\Sh(Y_{\et},n)\longrightarrow\Sh(Y_{\et},n)
\]
is left exact, and so we can form its right derived functors $\shExt^r(\F_0,-)$. These are called the local Ext groups. They are related to the global Ext groups by a spectral sequence
\[
H^r(Y_{\et},\shExt^s(\F_0,\G))\ \Longrightarrow\ \Ext^s(\F_0,\G).
\]
Under the hypothesis on $\F_0$, the stalk
\[
\shExt^r(\F_0,\G)_{\bar y}=\Ext^r((\F_0)_{\bar y},\G_{\bar y})
\]
for all geometric points $\bar y$. The Ext at right is formed in the category of $\Z/n\Z$-modules. Therefore, in the situation of the lemma
\[
\shExt^r(\F,\mu_n)_{\bar y}=\Ext^r(\F_{\bar y},\mu_n)
\approx\Ext^r(\F_{\bar y},\Z/n\Z)=0\quad\text{for }r>0,
\]
and the spectral sequence gives isomorphisms
\[
H^r(X_{\et},\shHom(\F,\mu_n))\simeq\Ext^r_{X,n}(\F,\mu_n),
\]
as required. For more details, see EC II 3.20, III 1.22, III 1.31.
\end{proof}
Thus, for a locally constant sheaf $\F$, the pairing in Theorem 14.20 takes the form
\[
H_c^r(U,\F)\times H^{2-r}(U,\check{\F}(1))
\longrightarrow H_c^2(U,\mu_n)\simeq\Z/n\Z.
\]

\loc{The pairing, p.\ 103; ending before ``The proof of Theorem 14.20''.}
There are various ways of defining the pairings, but (fortunately) they all agree, at least up to sign. Perhaps the simplest is to interpret the Ext groups as equivalence classes of extensions, and then define the pairing as an iterated boundary map.

In any abelian category, the elements of $\Ext^r(A,B)$, $r>0$, can be interpreted as equivalence classes of $r$-fold extensions
\[
0\longrightarrow B\longrightarrow E_1\longrightarrow\cdots
\longrightarrow E_r\longrightarrow A\longrightarrow0.
\]
There is a natural pairing
\[
\Ext^r(A,B)\times\Ext^s(B,C)\longrightarrow\Ext^{r+s}(A,C)
\]
that can either be defined by splicing extensions, or as an iterated boundary map, i.e., break the extension
\[
0\longrightarrow C\longrightarrow E_1\longrightarrow\cdots
\longrightarrow E_s\longrightarrow B\longrightarrow0
\]
into short exact sequences
\[
\begin{array}{rrrrrrrrrrr}
0&\longrightarrow&K_s&\longrightarrow&E_s&\longrightarrow&B&\longrightarrow&0\\
0&\longrightarrow&K_{s-1}&\longrightarrow&E_{s-1}&\longrightarrow&K_s&\longrightarrow&0\\
&&&&\vdots
\end{array}
\]
and form the boundary maps
\[
\Ext^r(A,B)\longrightarrow\Ext^{r+1}(A,K_s)
\longrightarrow\Ext^{r+2}(A,K_{s-1})
\longrightarrow\cdots\longrightarrow\Ext^{r+s}(A,C).
\]
See Mitchell, B., \textit{Theory of Categories}, Academic Press, 1965, VII.3 or Bourbaki, N., \textit{Alg\`ebre}, Chap.\ X.

In our case, an element of $\Ext^s(\F,\mu_n)$ is represented by an exact sequence
\[
0\longrightarrow\mu_n\longrightarrow E_1\longrightarrow\cdots
\longrightarrow E_s\longrightarrow\F\longrightarrow0.
\]
Apply $j_!$ to get an extension
\[
0\longrightarrow j_!\mu_n\longrightarrow j_!E_1\longrightarrow\cdots
\longrightarrow j_!E_s\longrightarrow j_!\F\longrightarrow0,
\]
which defines an iterated boundary map
\[
H^r(X,j_!\F)\longrightarrow H^{r+s}(X,j_!\mu_n).
\]
% END PAIRING CONTEXT LEC:14.7
\end{document}
